% ========== 基础包 ==========
\usepackage[UTF8]{ctex}               % 中文支持
\usepackage[margin=2.5cm]{geometry}   % 页边距
\usepackage{amsmath, amssymb}         % 数学公式
\usepackage{xcolor}                   % 颜色
\usepackage{mdframed}                 % 高亮框
\usepackage{listings}                 % 代码块
\usepackage{enumitem}                 % 列表样式
\usepackage{titlesec}                 % 标题样式
\usepackage{fancyhdr}                 % 页眉页脚
\usepackage{tocloft}                  % 目录样式
\usepackage{parskip}                  % 段落间距
\usepackage{graphicx}                 % 图片
\usepackage{pgfplots}                 % 函数图像
\pgfplotsset{compat=1.18}
\usepackage{tikz}                     % 绘图
\usetikzlibrary{arrows.meta, positioning, fit}  % tikz 库（bagel.tex 架构图需要）
\usepackage{array}                    % 高级表格列格式（>{\bfseries} 等）
\usepackage{booktabs}                 % 三线表（\toprule / \midrule / \bottomrule）
\usepackage{longtable}                % 跨页表格

% ========== 资源路径 ==========
\graphicspath{{assets/figures/}}

% ========== 颜色定义 ==========
\definecolor{primary}{RGB}{52, 120, 246}      % 主色：蓝
\definecolor{accent}{RGB}{255, 149, 0}        % 强调：橙
\definecolor{success}{RGB}{52, 199, 89}       % 成功：绿
\definecolor{danger}{RGB}{255, 59, 48}        % 警告：红
\definecolor{codebg}{RGB}{246, 248, 250}      % 代码背景
\definecolor{codeborder}{RGB}{208, 215, 222}  % 代码边框
\definecolor{blockbg}{RGB}{240, 247, 255}     % 信息框背景
\definecolor{textgray}{RGB}{100, 100, 100}    % 灰色文字

% ========== 标题样式 ==========
\titleformat{\section}
  {\Large\bfseries\color{primary}}
  {\thesection}{0.8em}{}
  [\vspace{-0.3em}\textcolor{primary}{\rule{\linewidth}{1.5pt}}\vspace{0.3em}]

\titleformat{\subsection}
  {\large\bfseries\color{black!80}}
  {\thesubsection}{0.8em}{}

\titleformat{\subsubsection}
  {\normalsize\bfseries\color{textgray}}
  {\thesubsubsection}{0.8em}{}

% ========== 页眉页脚 ==========
\pagestyle{fancy}
\fancyhf{}
\fancyhead[L]{\small\textcolor{textgray}{\leftmark}}
\fancyhead[R]{\small\textcolor{textgray}{KillerWang 的学习笔记}}
\fancyfoot[C]{\small\textcolor{textgray}{\thepage}}
\renewcommand{\headrulewidth}{0.4pt}
\renewcommand{\headrule}{\hbox to\headwidth{\color{codeborder}\leaders\hrule height \headrulewidth\hfill}}

% ========== 代码块样式 ==========
\lstset{
  backgroundcolor=\color{codebg},
  basicstyle=\ttfamily\small,
  keywordstyle=\color{primary}\bfseries,
  commentstyle=\color{textgray}\itshape,
  stringstyle=\color{success},
  frame=single,
  framerule=0.8pt,
  rulecolor=\color{codeborder},
  breaklines=true,
  breakatwhitespace=true,
  showstringspaces=false,
  tabsize=4,
  xleftmargin=8pt,
  xrightmargin=8pt,
  aboveskip=8pt,
  belowskip=8pt,
  numbers=left,
  numberstyle=\tiny\color{textgray},
  numbersep=8pt,
}

% ========== 自定义模块化环境 ==========

% 【关键概念框】蓝色
\newmdenv[
  backgroundcolor=blockbg,
  linecolor=primary,
  linewidth=2pt,
  leftline=true, rightline=false, topline=false, bottomline=false,
  innerleftmargin=12pt,
  innerrightmargin=8pt,
  innertopmargin=8pt,
  innerbottommargin=8pt,
  skipabove=8pt,
  skipbelow=8pt,
]{conceptbox}

% 【直觉理解框】橙色
\newmdenv[
  backgroundcolor={orange!8},
  linecolor=accent,
  linewidth=2pt,
  leftline=true, rightline=false, topline=false, bottomline=false,
  innerleftmargin=12pt,
  innerrightmargin=8pt,
  innertopmargin=8pt,
  innerbottommargin=8pt,
  skipabove=8pt,
  skipbelow=8pt,
]{intubox}

% 【警告/注意框】红色
\newmdenv[
  backgroundcolor={red!5},
  linecolor=danger,
  linewidth=2pt,
  leftline=true, rightline=false, topline=false, bottomline=false,
  innerleftmargin=12pt,
  innerrightmargin=8pt,
  innertopmargin=8pt,
  innerbottommargin=8pt,
  skipabove=8pt,
  skipbelow=8pt,
]{warnbox}

% 【总结框】绿色
\newmdenv[
  backgroundcolor={green!6},
  linecolor=success,
  linewidth=2pt,
  leftline=true, rightline=false, topline=false, bottomline=false,
  innerleftmargin=12pt,
  innerrightmargin=8pt,
  innertopmargin=8pt,
  innerbottommargin=8pt,
  skipabove=8pt,
  skipbelow=8pt,
]{summarybox}

% ========== 快捷命令 ==========
\newcommand{\key}[1]{\textbf{\textcolor{primary}{#1}}}        % 关键词
\newcommand{\warn}[1]{\textbf{\textcolor{danger}{#1}}}         % 警告
\newcommand{\note}[1]{\textit{\textcolor{textgray}{#1}}}       % 备注
\newcommand{\code}[1]{\texttt{\colorbox{codebg}{#1}}}          % 行内代码
\newcommand{\good}[1]{\textcolor{success}{\textbf{#1}}}        % 好的
\newcommand{\bad}[1]{\textcolor{danger}{\textbf{#1}}}          % 坏的

% ========== 超链接样式 ==========
\usepackage{hyperref}                 % 超链接
\hypersetup{
  colorlinks=true,
  linkcolor=primary,
  urlcolor=primary,
  citecolor=textgray,
}

% ========== 目录样式 ==========
\renewcommand{\cftsecfont}{\bfseries\color{primary}}
\renewcommand{\cftsecpagefont}{\bfseries\color{primary}}
